% コマ6: while / for の制御フローと break / continue の飛び先
\documentclass[border=8pt]{standalone}
% 講義資料の図(materials/sessions/figures/*.tex)共通プリアンブル
% 各図の .tex から % 講義資料の図(materials/sessions/figures/*.tex)共通プリアンブル
% 各図の .tex から % 講義資料の図(materials/sessions/figures/*.tex)共通プリアンブル
% 各図の .tex から \input{figure-preamble} で読み込む。
% 配色・フォントは latex/session-template.tex と揃える。

% 日本語(LuaLaTeX + luatexja)。等幅セル内の日本語も和文フォントで出す
\usepackage{luatexja}
\usepackage{luatexja-fontspec}
\setmainfont{Noto Sans CJK JP}
\setmainjfont{Noto Sans CJK JP}
\setmonofont{Inconsolata}[Scale=MatchLowercase]
\setmonojfont{Noto Sans CJK JP}

\usepackage{xcolor}
\definecolor{TcAccentDark}{HTML}{583B66}
\definecolor{TcAccentLight}{HTML}{EEE6F2}
\definecolor{TcAccentLighter}{HTML}{F7F3F9}
\definecolor{TcEmph}{HTML}{E64B17}
\definecolor{TcText}{HTML}{262626}
\definecolor{TcGrayText}{HTML}{737373}
\definecolor{TcGrayBg}{HTML}{F2F2F2}

\usepackage{tikz}
\usetikzlibrary{arrows.meta,positioning,calc,decorations.pathreplacing,fit,backgrounds}

\tikzset{
  % テキスト
  every node/.style={text=TcText},
  annot/.style={font=\small, text=TcText, align=left},
  gannot/.style={font=\small, text=TcGrayText, align=center},
  addr/.style={font=\ttfamily\small, text=TcGrayText},
  % メモリセル
  memcell/.style={draw=TcAccentDark, line width=0.6pt, fill=white,
                  minimum height=8.5mm, minimum width=40mm,
                  font=\ttfamily, inner sep=2pt, outer sep=0pt},
  memcell gray/.style={memcell, fill=TcGrayBg},
  memcell hi/.style={memcell, fill=TcAccentLighter},
  memcell dashed/.style={memcell, dashed, fill=none, text=TcGrayText},
  % 領域(セクション図など)
  region/.style={draw=TcAccentDark, line width=0.6pt, fill=white,
                 minimum width=48mm, align=center, inner sep=3pt, outer sep=0pt},
  % 制御フロー図のボックス
  cfgbox/.style={draw=TcAccentDark, line width=0.6pt, fill=white, rounded corners=1.5mm,
                 minimum width=28mm, minimum height=8mm, font=\small, align=center,
                 inner sep=2pt},
  % 矢印
  ptr/.style={-{Stealth[length=2.8mm]}, line width=1.1pt, color=TcEmph},
  flow/.style={-{Stealth[length=2.4mm]}, line width=0.8pt, color=TcAccentDark},
  flow back/.style={flow, dashed},
  regptr/.style={-{Stealth[length=2.8mm]}, line width=1.1pt, color=TcAccentDark},
  % 補助
  bracestyle/.style={decorate, decoration={brace, amplitude=4pt}, line width=0.8pt,
                     color=TcAccentDark},
}
 で読み込む。
% 配色・フォントは latex/session-template.tex と揃える。

% 日本語(LuaLaTeX + luatexja)。等幅セル内の日本語も和文フォントで出す
\usepackage{luatexja}
\usepackage{luatexja-fontspec}
\setmainfont{Noto Sans CJK JP}
\setmainjfont{Noto Sans CJK JP}
\setmonofont{Inconsolata}[Scale=MatchLowercase]
\setmonojfont{Noto Sans CJK JP}

\usepackage{xcolor}
\definecolor{TcAccentDark}{HTML}{583B66}
\definecolor{TcAccentLight}{HTML}{EEE6F2}
\definecolor{TcAccentLighter}{HTML}{F7F3F9}
\definecolor{TcEmph}{HTML}{E64B17}
\definecolor{TcText}{HTML}{262626}
\definecolor{TcGrayText}{HTML}{737373}
\definecolor{TcGrayBg}{HTML}{F2F2F2}

\usepackage{tikz}
\usetikzlibrary{arrows.meta,positioning,calc,decorations.pathreplacing,fit,backgrounds}

\tikzset{
  % テキスト
  every node/.style={text=TcText},
  annot/.style={font=\small, text=TcText, align=left},
  gannot/.style={font=\small, text=TcGrayText, align=center},
  addr/.style={font=\ttfamily\small, text=TcGrayText},
  % メモリセル
  memcell/.style={draw=TcAccentDark, line width=0.6pt, fill=white,
                  minimum height=8.5mm, minimum width=40mm,
                  font=\ttfamily, inner sep=2pt, outer sep=0pt},
  memcell gray/.style={memcell, fill=TcGrayBg},
  memcell hi/.style={memcell, fill=TcAccentLighter},
  memcell dashed/.style={memcell, dashed, fill=none, text=TcGrayText},
  % 領域(セクション図など)
  region/.style={draw=TcAccentDark, line width=0.6pt, fill=white,
                 minimum width=48mm, align=center, inner sep=3pt, outer sep=0pt},
  % 制御フロー図のボックス
  cfgbox/.style={draw=TcAccentDark, line width=0.6pt, fill=white, rounded corners=1.5mm,
                 minimum width=28mm, minimum height=8mm, font=\small, align=center,
                 inner sep=2pt},
  % 矢印
  ptr/.style={-{Stealth[length=2.8mm]}, line width=1.1pt, color=TcEmph},
  flow/.style={-{Stealth[length=2.4mm]}, line width=0.8pt, color=TcAccentDark},
  flow back/.style={flow, dashed},
  regptr/.style={-{Stealth[length=2.8mm]}, line width=1.1pt, color=TcAccentDark},
  % 補助
  bracestyle/.style={decorate, decoration={brace, amplitude=4pt}, line width=0.8pt,
                     color=TcAccentDark},
}
 で読み込む。
% 配色・フォントは latex/session-template.tex と揃える。

% 日本語(LuaLaTeX + luatexja)。等幅セル内の日本語も和文フォントで出す
\usepackage{luatexja}
\usepackage{luatexja-fontspec}
\setmainfont{Noto Sans CJK JP}
\setmainjfont{Noto Sans CJK JP}
\setmonofont{Inconsolata}[Scale=MatchLowercase]
\setmonojfont{Noto Sans CJK JP}

\usepackage{xcolor}
\definecolor{TcAccentDark}{HTML}{583B66}
\definecolor{TcAccentLight}{HTML}{EEE6F2}
\definecolor{TcAccentLighter}{HTML}{F7F3F9}
\definecolor{TcEmph}{HTML}{E64B17}
\definecolor{TcText}{HTML}{262626}
\definecolor{TcGrayText}{HTML}{737373}
\definecolor{TcGrayBg}{HTML}{F2F2F2}

\usepackage{tikz}
\usetikzlibrary{arrows.meta,positioning,calc,decorations.pathreplacing,fit,backgrounds}

\tikzset{
  % テキスト
  every node/.style={text=TcText},
  annot/.style={font=\small, text=TcText, align=left},
  gannot/.style={font=\small, text=TcGrayText, align=center},
  addr/.style={font=\ttfamily\small, text=TcGrayText},
  % メモリセル
  memcell/.style={draw=TcAccentDark, line width=0.6pt, fill=white,
                  minimum height=8.5mm, minimum width=40mm,
                  font=\ttfamily, inner sep=2pt, outer sep=0pt},
  memcell gray/.style={memcell, fill=TcGrayBg},
  memcell hi/.style={memcell, fill=TcAccentLighter},
  memcell dashed/.style={memcell, dashed, fill=none, text=TcGrayText},
  % 領域(セクション図など)
  region/.style={draw=TcAccentDark, line width=0.6pt, fill=white,
                 minimum width=48mm, align=center, inner sep=3pt, outer sep=0pt},
  % 制御フロー図のボックス
  cfgbox/.style={draw=TcAccentDark, line width=0.6pt, fill=white, rounded corners=1.5mm,
                 minimum width=28mm, minimum height=8mm, font=\small, align=center,
                 inner sep=2pt},
  % 矢印
  ptr/.style={-{Stealth[length=2.8mm]}, line width=1.1pt, color=TcEmph},
  flow/.style={-{Stealth[length=2.4mm]}, line width=0.8pt, color=TcAccentDark},
  flow back/.style={flow, dashed},
  regptr/.style={-{Stealth[length=2.8mm]}, line width=1.1pt, color=TcAccentDark},
  % 補助
  bracestyle/.style={decorate, decoration={brace, amplitude=4pt}, line width=0.8pt,
                     color=TcAccentDark},
}

\begin{document}
\begin{tikzpicture}[node distance=0mm,
    bc/.style={ptr, dashed}]
  % ================ while ================
  \node[gannot] at (0mm, 12mm) {{\ttfamily while (cond) \{ body \}}};

  \node[cfgbox] (Wc) at (0mm, 0mm) {cond を計算};
  \node[addr, left=1.5mm of Wc] {.Lcond:};
  \node[cfgbox, fill=TcAccentLighter] (Wb) at (0mm, -19mm) {body（本体の文）};
  \node[cfgbox] (We) at (0mm, -38mm) {以降の処理};
  \node[addr, left=1.5mm of We] {.Lend:};

  \draw[flow] (Wc) -- node[annot, right=1mm] {真} (Wb);
  \draw[flow] (Wb.west) to[out=180,in=180,looseness=1.6]
      node[annot, left=2mm, pos=0.5, align=right] {j .Lcond\\（繰り返す）}
      ($(Wc.west)+(0,-2mm)$);
  \draw[ptr] (Wc.east) to[out=0,in=0,looseness=1.6]
      node[annot, text=TcEmph, right=2mm, pos=0.5] {beqz a0, .Lend（偽）} (We.east);
  \draw[bc] ($(Wb.east)+(0,1.5mm)$) to[out=0,in=0,looseness=1.1]
      node[annot, text=TcEmph, right=1mm, pos=0.6] {continue\\（.Lcond へ）}
      ($(Wc.east)+(0,-2mm)$);
  \draw[bc] ($(Wb.east)+(0,-1.5mm)$) to[out=0,in=0,looseness=1.1]
      node[annot, text=TcEmph, right=1mm, pos=0.4] {break\\（.Lend へ）}
      ($(We.east)+(0,2mm)$);

  % ================ for ================
  \node[gannot] at (100mm, 12mm) {{\ttfamily for (init; cond; step) \{ body \}}};

  \node[cfgbox] (Fi) at (100mm, 0mm) {init を実行};
  \node[cfgbox] (Fc) at (100mm, -14mm) {cond を計算};
  \node[addr, left=1.5mm of Fc] {.Lcond:};
  \node[cfgbox, fill=TcAccentLighter] (Fb) at (100mm, -28mm) {body（本体の文）};
  \node[cfgbox] (Fs) at (100mm, -42mm) {step を実行};
  \node[addr, left=1.5mm of Fs] {.Lstep:};
  \node[cfgbox] (Fe) at (100mm, -56mm) {以降の処理};
  \node[addr, left=1.5mm of Fe] {.Lend:};

  \draw[flow] (Fi) -- (Fc);
  \draw[flow] (Fc) -- node[annot, right=1mm] {真} (Fb);
  \draw[flow] (Fb) -- (Fs);
  \draw[flow] ($(Fs.west)+(0,-2mm)$) to[out=180,in=180,looseness=1.2]
      node[annot, left=2mm, pos=0.5, align=right] {j .Lcond\\（繰り返す）}
      ($(Fc.west)+(0,-2mm)$);
  \draw[ptr] (Fc.east) to[out=0,in=0,looseness=1.7]
      node[annot, text=TcEmph, right=1.5mm, pos=0.85] {beqz a0, .Lend（偽）}
      (Fe.east);
  \draw[bc] ($(Fb.east)+(0,1.5mm)$) to[out=0,in=0,looseness=1.4]
      node[annot, text=TcEmph, right=1mm, pos=0.5] {continue（.Lstep へ！）}
      ($(Fs.east)+(0,1.5mm)$);
  \draw[bc] ($(Fb.west)+(0,-1.5mm)$) to[out=180,in=180,looseness=2.2]
      node[annot, text=TcEmph, left=1mm, pos=0.5, align=right] {break\\（.Lend へ）}
      ($(Fe.west)+(0,2mm)$);

  \node[annot, text=TcEmph, anchor=north, align=center] at (100mm, -66mm)
      {continue が .Lcond に飛ぶと、step が実行されず無限ループになる};

  % ================ 凡例 ================
  \draw[flow] (-16mm, -50mm) -- (-8mm, -50mm);
  \node[annot, anchor=west] at (-7mm, -50mm) {そのまま次の命令へ進む};
  \draw[ptr] (-16mm, -55mm) -- (-8mm, -55mm);
  \node[annot, anchor=west] at (-7mm, -55mm) {ジャンプ命令で飛ぶ};
  \draw[bc] (-16mm, -60mm) -- (-8mm, -60mm);
  \node[annot, anchor=west] at (-7mm, -60mm) {break / continue のジャンプ};
\end{tikzpicture}
\end{document}
